% Nouveau EGA V (cote 25): archive cote 25, site page 19 (PDF page 20).
% typescript. Draft compared: none.
% AI-transcribed by Claude Opus 5.5 (claude-opus-5-5, adaptive thinking, effort high), Message Batches API, batch ega25b (msgbatch_01A7n5YuXubUPNK1zsepi8RG), 2026-11-13; not yet checked by eye.
\scanpage{25/19}
\mspage{16}

sections hyperplanes de X lui-même (sans extension préalable du corps
\margin{\unsure{trait vertical et flèche}}de base). Cela résulte \unsure{en e ffet|en effet} de la constructibilité de E et du fait
que $\check{P}$ est un schéma de Jacobson.

Donnons à titre d'exemples quelques cas particuliers où les
considérations précédentes s'appliquent:

\uline{Proposition} 4.2. Soit G un Module de présentation finie sur X .
Soit P une des propriétés suivantes pour un Module\inserted{M} sur un préschéma
algébrique Z sur un corps K ;
\xout{\unsure{xxxthérixxxxxxx(i)xxxxxx(v)xxxxxxxxx2,8xxxAlxxxx}}
\deleted{a) L'ensemble \xout{des} E des \unsure{$\check{P}$} tels que G satisfasse P est
constructible}
(i) coppof (M) $\unsure{\leq}$ n.

(ii) M satisfait (S$_{k}$).

(i\unsure{\xout{x}}i) M est Cohen-Macaulay
\inserted[typed]{cycles}
(iv) M est sans \xout{idéaux} premiers associés immergés.

(v) M est séparable sur \xout{\illegible{1}} K .

\deleted{(vi) M est géométriquement intègre sur K/}

\xout{(vii)xxMxestxgéomxtxxxintègrexsurxKx} Avec ces notations; si E désig%
ne l'ensemble des \hand{$\xi\in$} P tels que G\hand{$_{\xi}$} satisfasse à P, on a :
a) E est une partie constructible de \over{P}{$\check{P}$} .

b) Supposons S=Spec(k), k un corps, et que F satisfasse à P. Suppo-%
sons de plus \inserted{(dans le cas (v))} que k soit de car 0, ou que f:X$\rightarrow$P soit une \deleted{immersion}\inserted{\unsure{non ramifié}}. \xout{\unsure{Akx}}
Alors E contient un ouvert dense de \over{P}{$\check{P}$} .

Démonstration. a) résulte du fait que P est une propriété
constructible, ce qui a été vu au par.9. Quant à b), cela résulte
des résultats correspondants des \xout{deux}\inserted{deux} n$^{o}$s précédents. Remords\inserted{à b)}: suppo-%
ser plus généralement que si Z est l'ensemble des points de X enlesquel%
F ne satisfait pas à P, on ait f(Z) fini i.e. dim $\overline{f(Z)}\;\unsure{\leq}0$.

\mspage{19}
